\chapter{Conclusion}
\label{ch:conclusion}

This project studied Affine Jump-Diffusion (AJD) models, with a focus on how their transform
structure gives both efficient pricing and calibration to market option data. 
By leaning on \parencite{Duffie2000}, we saw that the extended transform is exponential-affine in
the state variable, with coefficients given by a system of ODEs. Combined with the
\parencite{Lewis2001} formula, this gives European option prices as a single integral of the
characteristic function. For the SVJJ model, these ODEs admit closed-form solutions, hence no
numerical ODE solver is needed. 

We then calibrated SVJJ, SVJ and Heston to 1,290 out-of-the-money SPX option quotes across 12
maturities on 13 November 2019, after filtering for no-arbitrage violations. This analysis led to
three main findings:

\paragraph{The weighting scheme matters.}
The Vega-weighted objective gave the best overall fit in implied volatility for all three models,
and it was the only scheme under which no parameters hit their bounds for any model. We therefore
used it as the primary objective. Under Vega weighting, SVJJ attains an implied-volatility RMSE of
0.953 volatility points, against 0.955 for SVJ and 0.993 for Heston. The mean absolute error shows a
larger gap (0.452, 0.467 and 0.578, respectively).

\paragraph{Model complexity improves fit, but not uniformly.}
Adding jumps to the log-price improved the fit relative to Heston, most visibly in the left wing at
short maturities, where Heston underestimates the skew. The extra variance-jump component
of SVJJ gives only a marginal improvement over SVJ under Vega weighting, and under bid-ask weighting
SVJJ fits worse than both simpler models, with two parameters at their bounds. More flexible models
do not automatically calibrate better. This is likely a sign of a difficult objective landscape in
which the optimiser does not always find a minimum.

\paragraph{All three models struggle in the wings.}
They capture the general shape of the smile and fit the region around the money well, but they
deviate substantially in the far wings. At the longest maturity, all three understate the implied
volatility of deep OTM options.

\paragraph{Limitations \&  Extensions.}
Our analysis considered a single trading day, so the results may not carry over to other market
regimes. Even though we used 20 random starting points during calibration, we cannot rule out that
better local minima exist. Furthermore, the comparison is in-sample, so some of the flexibility of
SVJJ may reflect overfitting rather than genuine explanatory power. These limitations suggest
natural extensions. One is to calibrate across multiple dates, or more data to study parameter
stability and out-of-sample performance. Another test is to calibrate richer AJD specifications.
